\clearpage\thispagestyle{empty}
\pdfbookmark[0]{Quellen zum Kometenbericht vom Mai 1804}{comet180405-sources}
\begin{center}\textbf{Quellen zum Kometenbericht vom Mai 1804}\end{center}
\noindent Diese KI-Übertragung in heutiges Deutsch folgt dem Auszug aus dem Bericht von Franz Xaver von Zach in \href{https://gdz.sub.uni-goettingen.de/id/PPN236007467}{Gauß’ \textit{Werken}, Band~VI} (Göttingen, 1874), S.~247--248. Der Bericht erschien als Beitrag LIV im Maiheft 1804 der \textit{Monatlichen Correspondenz}, Band~IX, \href{https://www.digitale-sammlungen.de/en/view/bsb10538601?page=468}{S.~432}, \href{https://www.digitale-sammlungen.de/en/view/bsb10538601?page=469}{S.~433}, \href{https://www.digitale-sammlungen.de/en/view/bsb10538601?page=470}{S.~434} und \href{https://www.digitale-sammlungen.de/en/view/bsb10538601?page=471}{S.~435}. Von Zach berichtet über Beobachtungen von Heinrich Wilhelm Matthias Olbers, Jean-Louis Pons, Alexis Bouvard und Charles Messier und gibt die von Carl Friedrich Gauß bestimmten parabolischen Elemente und seine Berechnung wieder. Das auf S.~385 angeführte Verfahren stammt aus Gauß’ \href{https://www.digitale-sammlungen.de/en/view/bsb10538601?page=421}{\textit{Einige Bemerkungen zur Vereinfachung der Rechnung für die geocentrischen Oerter der Planeten}}, auch in den \textit{Werken}, Band~VI, S.~94.
\par\smallskip\noindent Übertragung und Satz wurden mit KI erstellt. Die Digitalisate der Zeitschrift stammen von der Bayerischen Staatsbibliothek.
\clearpage
